% !TeX root = ../thuthesis-example.tex

\chapter{系统设计与方法}

\section{总体架构设计}

本项目以 LingoBridge 现有国际中文学习平台为基础，将课程、作业、录音练习、批改反馈、学习时长、完成率和班级排名等业务数据转化为人工智能模块可处理的学习行为数据。系统整体采用“平台业务层 + 数据处理层 + 算法分析层 + 推荐展示层”的结构：平台业务层负责保留教师端、学生端和管理端的 MVP 主流程；数据处理层负责清洗学习日志并构造特征；算法分析层负责 DKT 知识追踪、Levenshtein 距离校准和知识点依赖分析；推荐展示层将结果反馈到学生首页、个人排名页、教师学生总览页和管理端数据中台。

该设计的核心目标不是重新开发一个独立算法 Demo，而是把人工智能方法嵌入真实教学平台，使系统既满足课程大作业对数据预处理、算法设计、交互设计和编码实现的要求，又能在本地演示环境中稳定运行。

\section{学习行为数据建模}

\subsection{数据来源与字段设计}

系统将学生学习过程抽象为时间序列日志。每条日志至少包含学生编号、知识点编号、练习类型、提交时间、答题结果、教师反馈、录音或字幕校准状态等字段。对于中文学习场景，知识点可以进一步划分为声母、韵母、声调、汉字、词汇和听读任务，从而形成适合知识追踪模型输入的细粒度标签。

\begin{equation}
r_t = (u, k_t, q_t, c_t, time_t, feedback_t)
\end{equation}

其中，$u$ 表示学生，$k_t$ 表示第 $t$ 次交互对应的知识点，$q_t$ 表示练习类型，$c_t$ 表示答题或任务完成结果，$time_t$ 表示发生时间，$feedback_t$ 表示教师或系统反馈。

\subsection{数据清洗与特征构建}

为了降低噪声对模型训练和推荐结果的影响，系统在进入算法模块前执行基础数据清洗：去除重复提交记录，补齐缺失知识点标签，统一时间格式，并按学生和知识点聚合练习次数、正确率、最近提交时间、连续错误次数和进步趋势。清洗后的特征既可用于 DKT 序列模型，也可作为管理端数据中台的统计指标。

\section{DKT 知识追踪模型设计}

DKT（Deep Knowledge Tracing）用于估计学生对不同中文知识点的动态掌握状态。系统将学生的历史练习序列编码为 $(k_t, c_t)$ 的输入序列，并通过循环神经网络学习知识状态随时间变化的规律。模型输出为学生在下一次练习中掌握各知识点的概率。

\begin{equation}
h_t = LSTM(x_t, h_{t-1})
\end{equation}

\begin{equation}
\hat{y}_{t+1} = \sigma(W h_t + b)
\end{equation}

其中，$x_t$ 是第 $t$ 次学习交互的编码向量，$h_t$ 是隐状态，$\hat{y}_{t+1}$ 表示模型预测的下一时刻知识点掌握概率。对于本地课程展示，系统采用轻量 LSTM 原型，使模型可以在普通 CPU 环境下完成训练和推理。

\section{Levenshtein 距离校准方法}

在中文学习中，拼音、声调和字幕转写常出现相近发音导致的偏误。系统采用 Levenshtein 编辑距离对 ASR 字幕、拼音转写或练习答案进行相似度校准。设原始识别文本为 $s$，标准候选文本为 $w$，系统在词表中搜索编辑距离最小的候选项：

\begin{equation}
w^* = \arg\min_{w \in Dict} Levenshtein(s, w)
\end{equation}

当最小距离低于阈值时，系统认为该文本存在可校准偏误，并将校准结果作为推荐理由的一部分。例如，若学生在拼音或字幕中反复出现相近声调错误，系统会把对应声调知识点标记为薄弱项，并推荐下一轮专项练习。

\section{知识点依赖与推荐策略}

除了 DKT 和 Levenshtein 校准外，系统还使用知识点依赖图表达声母、韵母、声调、汉字和词汇之间的先后关系。该图结构用于解释推荐路径：当学生在某一复合知识点上表现较弱时，系统可以回溯到更基础的声母、韵母或声调节点，避免只给出表层题目推荐。

最终推荐分数由知识掌握概率、最近错误情况、校准偏误次数和知识点依赖权重共同决定：

\begin{equation}
Score(k) = \alpha (1 - Mastery(k)) + \beta Error(k) + \gamma Calib(k) + \delta Dep(k)
\end{equation}

其中，$Mastery(k)$ 表示 DKT 输出的掌握概率，$Error(k)$ 表示历史错误强度，$Calib(k)$ 表示 Levenshtein 校准发现的偏误强度，$Dep(k)$ 表示知识点依赖权重。系统按照分数排序生成个性化练习推荐，并在学生端和教师端用可解释文案展示推荐依据。
